\chapter{Mathematical Background}

\section{Vectors and Matrices}

These notes will skip some preliminaries that you should know from first and second year. \\

If \(\fx^T y = 0\), then we say the two vectors are \(\fx\) and \(\fy\) are orthogonal.

\defn{Vector Norm}{
    A \textit{vector norm} on \(\bR^n\) is a function \(\norm{\cdot}\) from \(\bR^n\) to \(\bR\) such that
    \begin{enumerate}
        \item \(\norm{\fx} \geq 0\) for all \(\fx \in \bR^n\) and \(\norm{\fx} = 0 \iff \fx = 0\).
        \item \(\norm{\fx + \fy} \leq \norm{\fx} + \norm{\fy}\) for all \(\fx, \fy \in \bR^n\). (Triangle inequality)
        \item \(\norm{\alpha \fx} = |\alpha|\norm{\fx}\) for all \(\alpha \in \bR, \fx \in \bR^n\).
    \end{enumerate}
}

An \((n \times n)\) matrix \(A\) is said to be symmetric if \(A = A^T\). We consider eigenvalues in non-ascending order. \\

\textbf{Theorem. (Spectral Decomposition Theorem)}
Let \(A \in \bR^{n \times n}\) be a symmetric \(n \times n\) matrix. Then, there exists an orthogonal matrix \(Q\) (that is, \(Q^TQ = QQ^T = I_n\)) such that \(A = QDQ^T\) where \(D = \diag(\lambda_1, \dots, \lambda_n)\) and \(Q = \begin{bmatrix}
    \fv_1 & \fv_2 & \cdots & \fv_n
\end{bmatrix}\) where \(\fv_i\) is an eigenvector of \(A\) corresponding to the eigenvalue \(\lambda_i\). \\

For a symmetric \((n \times n)\) matrix \(A\) and any \(\fx \neq 0\), we define the so-called Rayleigh quotient as
\[R_A(\fx) = \frac{\fx^T A\fx}{\norm{\fx}_2^2}.\]
Then, we have, for any \(\fx \neq \fzero\),
\[\lambda_{\min}(A) \leq R_A(\fx) \leq \lambda_{\max}(A).\]

\textbf{Theorem. (Variational Characterization for Extremal Eigenvalues)}
For a symmetric \((n \times n)\) matrix \(A\),
\[\lambda_{\max}(A) = \max\{R_A(\fx) : \fx \neq 0\} = \max_{\fx \in \bR^n}\{\fx^T A\fx : \norm{\fx}_2 = 1\}\]
and
\[\lambda_{\min}(A) = \min\{R_A(\fx) : \fx \neq 0\} = \min_{\fx \in \bR^n} \{\fx^T A\fx : \norm{fx}_2 = 1\}.\]

\defn{Positive Definite Matrices}{
    A matrix \(A \in \bR^{n \times n}\) is
    \begin{itemize}
        \item \textit{positive definite} \(\iff \fx^T A\fx > 0\) for all \(\fx \in \bR^n, \fx \neq 0\)
        \item \textit{positive semi-definite} \(\iff \fx^T A\fx \geq 0\) for all \(\fx \in \bR^n\)
        \item \textit{negative definite} \(\iff \fx^T A\fx \leq 0\) for all \(\fx \in \bR^n\)
        \item \textit{indefinite} \(\iff\) there exists \(\fx, \fy \in \bR^n : \fx^T A\fx > 0\) and \(\fy^T A \fy < 0\)
    \end{itemize}
}

A symmetric matrix \(A \in \bR^{n \times n}\)
\begin{itemize}
    \item has \(n\) real eigenvalues \(\lambda_i, i = 1, \dots, n\)
    \item Determinant: \(\det(A) = \prod_{i=1}^n \lambda_i\)
    \item Trace: \(\trace(A) := \sum_{i = 1}^n a_{ii} = \sum_{i = 1}^n \lambda_i\)
\end{itemize}

\prop{
    A symmetric matrix \(A \in \bR^{n \times n}\) is
    \begin{itemize}
        \item \textit{positive definite} \(\iff \lambda_i > 0\) for all \(i = 1, \dots, n\)
        \item \textit{positive semi-definite} \(\iff \lambda_i \geq 0\) for all \(i = 1, \dots, n\)
        \item \textit{negative definite} \(\iff \lambda_i < 0\) for all \(i = 1, \dots, n\)
        \item \textit{negative semi-definite} \(\iff \lambda_i \leq 0\) for all \(i = 1, \dots, n\)
        \item \textit{indefinite} \(\iff\) there exists \(i, j: \lambda_i > 0\) and \(\lambda_j < 0\)
    \end{itemize}
}

\thrm{Checking Positive (Semi-)definitness}{
For any eigenvalues \(\lambda\) of a symmetric matrix \(A \in \bR^{n \times n}\), there exists \(i_0 \in \{1, \dots, n\}\) such that
\[\lambda - A_{i_{0}i_{0}} \leq \sum_{j \neq i_0} |A_{i_0}j.\]
}

\defn{Diagonally Dominant}{
A symmetric matrix \(A \in \bR^{n \times n}\) is called diagonally dominant if
\[|A_{ii}| \geq \sum_{j \neq i} |A_{ij}|, \text{ for each } i \in \{1, \dots, n\}.\]
In addition, if
\[|A_{ii}| > \sum_{j \neq i}|A_{ij}|, \text{ for each } i \in \{1, \dots, n\},\]
then we say \(A\) is strictly diagonally dominant.
}

Let \(A \in \bR^{n \times n}\) be a symmetric matrix. Suppose that \(A\) is diagonally dominant and all of its diagonal elements are non-negative, then any eigenvalue of \(A\) must be non-negative, and so, \(A\) is positive semi-definite.

\defn{Matrix Norm}{
    A \textit{matrix norm} on \(\bR^{p \times n}\) is a function \(\norm{\cdot}\) from \(\bR^{p \times n}\) to \(\bR\) such that
    \begin{enumerate}
        \item \(\norm{X} \geq 0\) for all \(X \in \bR^{p \times n}\) and \(\norm{X} = 0 \iff X = 0_{p \times n}\).
        \item \(\norm{X + Y} \leq \norm{X} + \norm{Y}\) for all \(X, Y \in \bR^{p \times n}\). (Triangle inequality)
        \item \(\norm{\alpha X} = |\alpha|\norm{X}\) for all \(\alpha \in \bR, X \in \bR^{p \times n}\).
    \end{enumerate}
}

One of the most widely used matrix norm is: \textbf{Frobenius norm}: \(\norm{X}_F = \left(\sum_{i = 1}^p \sum_{j = 1}^n X_{ij}^2\right)^\frac{1}{2}\). \\

We can also define a dot product between two matrices \(A, B \in \bR^{p \times n}\) as
\[A \cdot B = \trace(A^T B).\]

\section{Local and Global Minima}

Let \(H\) be a finite dimensional space and let \(\Omega\) be a set in \(H\). Let \(f: H \to \bR\) be a function on \(H\) and let \(\norm{\cdot}\) be a norm in \(H\).

\defn{Global Minimum}{
    A point \(\fx^* \in \Omega\) is a \textit{global minimizer} of \(f(\fx)\) over \(\Omega \subseteq H\) if and only if \(f(\fx^*) \leq f(\fx)\) for all \(\fx \in \Omega\). The global minimum is \(f(\fx^*)\).
}

\textbf{Definition. (Strict Global Minimum)} A point \(\fx^* \in \Omega\) is a \textit{strict global minimizer} of \(f(\fx)\) over \(\Omega \subseteq H\) if and only if \(f(\fx^*) < f(\fx)\) for all \(\fx \in \Omega, \fx \neq \fx^*\).

\defn{Local Minimum}{
    A point \(\fx^* \in \Omega\) is a \textit{local minimizer} of \(f(\fx)\) over \(\Omega \subseteq H\) if and only if there exists a \(\delta > 0\) such that \(f(\fx^*) \leq f(\fx)\) for all \(\fx \in \Omega\) with \(\norm{\fx - \fx^*} \leq \delta\). Then \(f(\fx^*)\) is a local minimum.
}

\textbf{Definition. (Strict Local Minimum)} A point \(\fx^* \in \Omega\) is a \textit{strict local minimizer} of \(f(\fx)\) over \(\Omega \subseteq H\) if and only if there exists a \(\delta > 0\) such that \(f(\fx^*) < f(\fx)\) for all \(\fx \in \Omega\) with \(0 < \norm{\fx - \fx^*} \leq \delta\).

\defn{Global Maximum}{
    A point \(\fx^* \in \Omega\) is a \textit{global maximizer} of \(f(\fx)\) over \(\Omega \subseteq H\) if and only if \(f(\fx^*) \geq f(\fx)\) for all \(\fx \in \Omega\). The global minimum is \(f(\fx^*)\).
}

\textbf{Definition. (Strict Global Maximum)} A point \(\fx^* \in \Omega\) is a \textit{strict global maximizer} of \(f(\fx)\) over \(\Omega \subseteq H\) if and only if \(f(\fx^*) > f(\fx)\) for all \(\fx \in \Omega, \fx \neq \fx^*\).

\defn{Local Maximum}{
    A point \(\fx^* \in \Omega\) is a \textit{local maximizer} of \(f(\fx)\) over \(\Omega \subseteq H\) if and only if there exists a \(\delta > 0\) such that \(f(\fx^*) \geq f(\fx)\) for all \(\fx \in \Omega\) with \(\norm{\fx - \fx^*} \leq \delta\). Then \(f(\fx^*)\) is a local minimum.
}

\textbf{Definition. (Strict Local Maximum)} A point \(\fx^* \in \Omega\) is a \textit{strict local maximizer} of \(f(\fx)\) over \(\Omega \subseteq H\) if and only if there exists a \(\delta > 0\) such that \(f(\fx^*) > f(\fx)\) for all \(\fx \in \Omega\) with \(0 < \norm{\fx - \fx^*} \leq \delta\).

\defn{Closed Sets}{
    A set \(\Omega\) is said to be closed if it contains all the limits of convergent sequences of points in \(\Omega\).
}

\defn{Bounded Sets}{
    A set \(\Omega\) is said to be bounded if there exists \(R > 0\) such that \(\Omega \subseteq \{x : \norm{x} \leq R\}\).
}

\textbf{Proposition (Existence of Global Solution: Type I)} Let \(H\) be a finite dimensional space and let \(\Omega\) be a closed and bounded set in \(H\) and let \(f\) be continuous on \(\Omega\). Then the global minima of \(f\)over \(\Omega\) exist. \\

\textbf{Proposition (Existence of Global Solution: Type II)} Let \(H\) be a finite dimensional space and let \(\norm{\cdot}\) be a norm in \(H\). Let \(\Omega\) be a closed set in \(H\) and let \(f\) be continuous on \(\Omega\). Suppose that \(f\) is coercive in the sense that

\[\lim_{\norm{x} \to \infty} f(x) = +\infty.\]

Then the global minima of \(f\) over \(\Omega\) exist. \\

\textbf{Relaxation} Let \(H\) be a finite dimensional space. If \(f : H \to \bR\) and \(\overline{\Omega} \subseteq \Omega\) then
\[\min_{x \in \Omega} f(\fx) \leq \min_{\fx \in \overline{\Omega}} f(\fx).\]